\documentclass[11pt, letterpaper]{article}
\usepackage[utf8]{inputenc}
\usepackage[T1]{fontenc}
\usepackage{geometry}
\geometry{letterpaper, margin=1.0in}
\usepackage{graphicx}
\usepackage{hyperref}
\usepackage{booktabs}
\usepackage{titlesec}
\usepackage{enumitem}
\usepackage{xcolor}
\usepackage{fancyhdr}
\usepackage{setspace}

% Professional Formatting
\setlength{\parindent}{0pt}
\setlength{\parskip}{1em}
\linespread{1.2}

% Section Styling
\titleformat{\section}
  {\Large\bfseries}
  {\thesection}{1em}{}[\titlerule]
  
\titleformat{\subsection}
  {\large\bfseries\itshape}
  {\thesubsection}{1em}{}

% Header/Footer
\pagestyle{fancy}
\fancyhf{}
\rhead{\textit{Future of Jobs - Project Report}}
\lhead{Abhishek M. Sutaria}
\cfoot{\thepage}

\title{\textbf{\huge Future of Jobs: Technical Realization \& Development Status}}
\author{
    \large \textbf{Abhishek Manoj Sutaria} \\ 
    Kelley School of Business \\[2em] 
    Prepared for: Professor Raymond Burke \\
    Date: February 2026
}
\date{}

\begin{document}

\maketitle
\thispagestyle{empty}

\section{System Architecture \& Improvements}

This report outlines the comprehensive transformation of the \textit{Future of Jobs} platform. We have moved beyond a simple concept to a fully functional, data-driven application. The following sections detail the completed improvements and the active development features currently underway.

\subsection{Data Integrity: From Assumption to Fact}
The core value of the platform has been secured by effectively replacing "simulated" data with verified, real-world sources.
\begin{enumerate}
    \item \textbf{U.S. Bureau of Labor Statistics (BLS):} We now directly integrate verified data for over 413,000 employment records. This means the specific salary numbers, growth projections, and location data are now factual and authoritative, rather than rough estimates.
    \item \textbf{Google Gemini 2.0 Integration:} Where historical data is too slow to catch up, we use advanced Artificial Intelligence to predict emerging trends, giving the platform a "living" analysis of the job market.
\end{enumerate}

\newpage
\subsection{Feature Expansion: Candidate Relevance Engine}
To make the tool personal for every student, we built a \textbf{Resume Analysis Module}.
\begin{itemize}
    \item \textbf{Instant PDF Scanning:} Students can drag and drop their actual PDF resumes. The system instantly reads them in the browser by parsing the pdf at the same instant and feeding it to AI.
    \item \textbf{AI "Survival Score":} The system doesn't just look for keywords; it understands the student's career path and compares it against AI automation trends to give a personalized "5-Year Relevance Score."
    \item \textbf{Actionable Advice:} Instead of generic feedback, it gives specific advice—for example, telling a student to "Shift focus from Manual Data Entry to Strategic Analytics" to stay competitive.
\end{itemize}

\subsection{User Interface (UI) Overhaul}
Based on direct feedback, the interface was completely redesigned to be intuitive:
\begin{itemize}
    \item \textbf{Visual Precision \& Alignment:} We meticulously aligned all buttons, cards, and data panels to a strict grid system. This ensures the interface looks clean, professional, and visually balanced, preventing any "cluttered" feeling.
    \item \textbf{Command Center:} The dashboard was simplified into a clean, 4-step journey, making powerful insights easy to access without digging through menus.
\end{itemize}

\section{Technical Engineering \& Optimization}
\textit{Creating a seamless experience required significant engineering effort "under the hood" to ensure the application is fast, stable, and professional.}

\begin{itemize}
    \item \textbf{Complete Code Reorganization (Stability):} The original code was massive and complex. I manually deconstructed it into smaller, organized pieces. This "modular" approach means the app is much more stable and new features can be added without breaking existing ones.
    \item \textbf{Zero-Error Standard (Reliability):} I conducted a strict review of the entire codebase, fixing multiple potential issues. This ensures the application doesn't crash or behave inextricably, providing a professional-grade experience.
    \item \textbf{Performance for Everyone (Optimization):} The 3D graphics are heavy, but I engineered an optimization system that automatically adjusts quality. This ensures the stunning visuals work smoothly on a standard student laptop, not just expensive high-end computers.
\end{itemize}

\section{User Evaluation \& Feedback}
We tested the platform with real users, including \textbf{Neel Shah}, \textbf{Aaryan Purohit}, and \textbf{Dev Patel}. Their feedback directly shaped our improvements:
\begin{itemize}
    \item \textbf{Clarity:} They found the original scores confusing, so we split them into clear "Automation Risk" vs. "Human Skills" metrics.
    \item \textbf{Next Steps:} They asked, "What do I do now?"—so we built the "Upskilling Plan" feature to give them a direct answer.
\end{itemize}

\section{Current Features in Development}

\subsection{Ontology Expansion: The "Kelley 50"}
I am currently working to expanding the job list from the initial 13 roles to the \textbf{Top 50 Marketing \& Business Roles} most relevant to Kelley students. This will ensure every student finds their specific dream job in the system.

\subsection{Dynamic Timeline Engine}
I am developing a "Rolling Timeline" feature. This means if a student uses the app in 2027, the system will automatically project 5 years into the future (to 2032). This ensures the tool never becomes outdated and always offers a relevant 5-year outlook.

\end{document}
